% 摘要
本文针对策略类游戏"代号K"的新服运营优化问题，基于2,000名玩家46天的真实行为数据，综合运用生存分析、机器学习、因果推断与蒙特卡洛模拟方法，建立了从留存规律挖掘到付费策略设计的完整模型体系，依次求解四个递进问题。

\textbf{问题1——留存预测。}严格限定前3天数据以避免未来信息泄漏。Kaplan-Meier估计给出次日留存率33.55\%、30日留存率仅0.85\%，中位生存时间1天。Cox比例风险模型取得C-index=0.899，分段分析揭示前3天付费用户的7日留存率是非付费用户的7.8倍，入盟用户是未入盟用户的23倍——付费与社交是留存的最强保护因子。

\textbf{问题2——付费关联。}使用全周期数据识别出矿石消耗与等级增长相关性最强，Logistic回归确定钻石存量流失风险阈值为566。等级停滞分析揭示4--6级为关键瓶颈，钻石是打破瓶颈的保护因子、金币是停滞堆积物。两阶段Hurdle模型进一步揭示付费概率由行为广度驱动、付费金额由资源消费强度驱动。

\textbf{问题3——策略设计。}K-Means聚类将玩家分为高氪核心党、中氪战力党、零氪休闲党和零氪流失党四种类型，构建基线—基准方案—Uplift优化方案三组递进对比体系。5,000次蒙特卡洛模拟验证优化方案实现71,045元期望营收，留存率11.3\%，$P(\geq 70,000)=52.4\%$，同时满足留存$\geq 10\%$约束与70,000元营收目标。

\textbf{问题4——数据采集。}建议优先采集社交行为与实时进度两类埋点数据，并设计AB测试闭环验证框架。

本文提出的"基准方案+Uplift优化+敏感性分析"范式，为无A/B测试条件下的游戏运营策略优化提供了可复用的方法论。

\vspace{1em}
\noindent\textbf{关键词：}生存分析；Cox比例风险模型；两阶段Hurdle模型；Uplift建模；蒙特卡洛模拟；付费策略优化
